\chapter[1910 --- Kossel]{1910: Albrecht Kossel}
\label{ch:1910_albrechtkossel}

\section*{Main Contribution}

\textit{Identified histidine as a fundamental building block of proteins, advancing understanding of protein chemistry and cellular metabolism.}

\section{Official Citation}

in recognition of the contributions to our knowledge of cell chemistry

\section{Background and Historical Context}

Albrecht Kossel was born on September 16, 1853, in Rostock, a port city on the Baltic coast of northern Germany. The son of a merchant family with roots in Pomerania, Kossel studied medicine at the Universities of Strasbourg, W\"urzburg, and Berlin, where he trained under some of the most distinguished biochemists and physiologists of the era, including Felix Hoppe-Seyler, the founder of the field of biochemistry. Kossel followed Hoppe-Seyler to the University of Strasbourg and remained there as a researcher and teacher for most of his career, eventually becoming professor of physiological chemistry and director of the Institute for the Study of Physiological Chemistry.

The late nineteenth century was a period of rapid advance in the understanding of the chemical composition of living cells. The new science of biochemistry---the chemistry of biological processes---was emerging as a distinct discipline, driven by the application of chemical methods to the study of biological materials. Proteins, fats, and carbohydrates were being isolated and characterized, and the first enzymes were being discovered. But the most fundamental question in biochemistry---What are the chemical constituents of the cell nucleus, and what is their function?---remained unanswered.

The cell nucleus had been recognized as a distinct organelle since the 1830s, when Robert Brown described it in plant cells. By the 1860s, it was well established that the nucleus was present in virtually all eukaryotic cells and that it played a central role in cell division. Friedrich Miescher, a Swiss biochemist working in T\"ubingen, had isolated a novel substance from the nuclei of white blood cells in 1869, which he called ``nuclein.'' This substance was rich in phosphorus and appeared to be a combination of a protein and an unknown phosphorus-containing compound. Miescher's nuclein was the first description of what we now call nucleic acid---the chemical substance of the genes. However, the significance of Miescher's discovery was not recognized for decades, and the chemical nature and function of nuclein remained mysterious.

\section{The Discovery}

Kossel's Nobel Prize-winning work was a systematic investigation of the chemical composition of cell proteins and, crucially, the nucleic substances. Through decades of meticulous biochemical research, he identified and characterized the fundamental building blocks of proteins and nucleic acids---work that laid the foundation for modern molecular biology.

\subsection{Identification of Amino Acids}

Kossel's first major contribution was the identification and characterization of several of the amino acids---the building blocks of proteins. While some amino acids had been isolated from natural sources before Kossel's work, he systematically identified and named the basic (alkaline) amino acids found in proteins. He isolated histidine (1896), arginine (1895), and lysine (1899) from protein digests, demonstrating that proteins were composed of distinct chemical subunits that could be separated and characterized. He also showed that these basic amino acids were particularly abundant in the histone proteins found in cell nuclei---an observation that hinted at the importance of these proteins in nuclear function.

Kossel's identification of amino acids was technically demanding work. Proteins are large, complex molecules that cannot be directly analyzed by the methods available to nineteenth-century chemists. Kossel developed techniques for the systematic hydrolysis (chemical breakdown) of proteins into their component amino acids, followed by separation and identification of the individual amino acids using precipitation, crystallization, and elemental analysis. His identification of the basic amino acids---histidine, arginine, and lysine---was recognized by the Nobel Committee as a major contribution to protein chemistry.

\subsection{Characterization of Nucleic Acids}

Kossel's most historically significant contribution was his characterization of the nucleic substances---the compounds that we now call DNA (deoxyribonucleic acid) and RNA (ribonucleic acid). Building on Miescher's discovery of nuclein, Kossel and his students isolated and characterized the components of nucleic acids.

Kossel demonstrated that nucleic acids were composed of two types of building blocks: (1) nitrogenous bases---adenine, guanine, cytosine, thymine (in DNA), and uracil (in RNA); and (2) a sugar component---deoxyribose (in DNA) or ribose (in RNA). He also showed that nucleic acids contained a phosphate group, and that the bases, sugar, and phosphate were linked together in a specific pattern to form nucleotides, the monomeric units of nucleic acids.

The identification of these components was achieved through painstaking chemical analysis. Kossel used acid hydrolysis to break nucleic acids into their component parts, and then separated and identified the products using a combination of precipitation, crystallization, and elemental analysis. His identification of the purine bases adenine and guanine, and the pyrimidine bases cytosine and thymine, established the chemical vocabulary that would later be used to describe the structure of DNA.

Kossel also made the important observation that the nucleic acids extracted from different organisms had different chemical compositions, suggesting that nucleic acids might carry organism-specific information---an insight that anticipated the discovery of DNA as the carrier of genetic information.

\subsection{Histones}

Kossel's identification of histones---a class of basic proteins found in cell nuclei---was another important contribution. He showed that histones were rich in the basic amino acids he had identified (lysine and arginine) and that they were closely associated with nucleic acids in the nucleus. The function of histones was not understood at the time, but subsequent research has shown that histones play a crucial role in packaging DNA and regulating gene expression. The histone code hypothesis---the idea that modifications to histone proteins regulate the accessibility and transcription of genes---is now a major area of research in molecular biology and epigenetics.

\section{Impact on Modern Medicine}

Kossel's characterization of the chemical components of nucleic acids laid the groundwork for one of a major scientific discoveries of the twentieth century: the determination of the structure of DNA by James Watson and Francis Crick in 1953. Watson and Crick's double-helix model of DNA, which revealed the molecular basis of genetic inheritance, was built directly on the chemical foundations that Kossel had established. The nucleotide bases, the sugar-phosphate backbone, and the complementary base pairing that are the structural features of DNA all depend on the chemical components that Kossel identified.

The understanding of DNA structure and function, which Kossel's work made possible, has transformed medicine. Genetic testing for inherited diseases, gene therapy for genetic disorders, pharmacogenomics (the tailoring of drug therapy to an individual's genetic makeup), and the molecular diagnosis of cancer through DNA sequencing are all direct descendants of the biochemical foundation that Kossel laid.

Kossel's work on amino acids also contributed to the understanding of protein chemistry that is essential for modern medicine. The identification of amino acids as the building blocks of proteins was a prerequisite for understanding protein structure, function, and folding---knowledge that is essential for the design of biologic drugs (including monoclonal antibodies, enzymes, and hormones) and for understanding the molecular basis of diseases caused by protein misfolding, such as Alzheimer's disease and sickle cell anemia.

The histones that Kossel identified are now recognized as fundamental components of the epigenetic regulatory machinery. Modifications to histone proteins---acetylation, methylation, phosphorylation, and ubiquitination---regulate the accessibility of genes to the transcriptional machinery and play crucial roles in development, differentiation, and disease. Epigenetic therapies, including histone deacetylase inhibitors, are now used in the treatment of certain cancers, and the study of the ``histone code'' is one of the most active areas of biomedical research.

\section{Legacy}

Albrecht Kossel received the Nobel Prize in Physiology or Medicine in 1910 ``in recognition of the contributions to our knowledge of cell chemistry made through his work on proteins, including the nucleic substances.'' He was the first scientist to receive the Nobel Prize for work in biochemistry, and his award recognized the growing importance of biochemical methods for understanding the fundamental processes of life.

Kossel continued his research in Strasbourg until his retirement in 1923. He died on July 5, 1927, at the age of 73. His laboratory in Strasbourg became a center for biochemical research and training, and his students carried his methods and ideas to institutions throughout Europe.

Kossel's legacy is the chemical foundation of molecular biology. By identifying the building blocks of proteins and nucleic acids, he provided the chemical vocabulary and analytical methods that made possible the subsequent revolution in molecular genetics. In an era when the sequencing of entire genomes, the editing of genes with CRISPR technology, and the development of personalized medicine based on an individual's genetic makeup are transforming medicine, the importance of Kossel's contribution to the chemical understanding of life can hardly be overstated.


\section{Modern Developments}

Kossel's work on nucleic acid chemistry led directly to the discovery of DNA structure. Modern genomics, including whole-genome sequencing, is built on understanding the nucleotide building blocks Kossel identified. The Human Genome Project (2003) and the Telomere-to-Telomere Consortium (2022) have completed the first gapless sequence of a human genome. RNA therapeutics, including mRNA vaccines (Nobel 2023), exploit the nucleotide chemistry Kossel first characterized.
